\subsection{Données et classes}
\label{sec:data}

\textbf{T3 (dépistage).}
Fusion de Kvasir v1 et HyperKvasir en trois classes (repère, muqueuse normale, polype).
Total : \numsamples{} images, centre \texttt{simula}.
Splits \textbf{groupés} par procédure/vidéo (\texttt{prepare-colonoscopy-3class} lit \texttt{image-labels.csv} HyperKvasir).
Les images publiques actuelles sont une frame par groupe ; les groupes multi-frames s'appliquent avec les pools vidéo non annotés.

\textbf{T1 (HyperKvasir-23).}
23 dossiers officiels Simula ; 10\,662 images, split 70/15/15.

Extensions prévues (hors tableaux principaux) : sous-ensemble pathologie-11 ; affinages JEPA sur HK-23.

\subsection{Protocole expérimental}
\label{sec:protocol}

\textbf{T3 (C0--C5, C7--C9).} Dépistage trois classes sur split \textbf{poolé} Simula (centre public unique ; pas de LOO multi-centres). C7--C9 varient la fraction de labels.

\textbf{T1 (B2--B6).} ResNet, SE-ResNet, dual-stream sur HyperKvasir-23 (supervisé uniquement dans ce manuscrit).

\subsection{Métriques}
Macro-F1 principale ; \textbf{rappel polype} $= \mathrm{TP}/(\mathrm{TP}+\mathrm{FN})$ pour la classe polype (sensibilité).

\subsection{Disponibilité des données et éthique}
\label{sec:ethics}

\textbf{Données.} Kvasir v1 et HyperKvasir sont des jeux publics Simula ; nous utilisons les dossiers officiels et \texttt{image-labels.csv} pour le regroupement par procédure.
Aucune donnée patient privée ni vidéo dans les runs rapportés.

\textbf{Éthique.} Images fixes dé-identifiées ; analyse secondaire de benchmarks ouverts, pas d'étude clinique prospective.
